\documentclass[11pt]{article}
\usepackage{amsmath,amssymb,amsthm}
\newtheorem{proposition}{Proposition}
\title{One good run, arbitrary bad complement, and the original sparse record}\author{}\date{}
\begin{document}\maketitle
\paragraph{Inputs and fixed-threshold clarification.}
The estimates of the path centered growing bath tubes calculation extend to a FIXED sufficiently small angular threshold
$\epsilon_0>0$: its centered Hessian differs from the flat one by
$O(\epsilon_0)+o_B(1)$, its tails have the same uniform Gaussian bound, and its
operator distance can therefore be made $<1/20$ by first choosing $\epsilon_0$
small and then $B$ large. Literal convergence to the flat kernel is not claimed
at fixed $\epsilon_0$. All subsequent cone constants remain unchanged.
We use only this fixed-threshold version. The exact bath criterion in the original sparse bath records calculation
puts a good interval into its required frame by left/right Haar transport;
all endpoint rotations are retained.

In the original sparse bath records calculation, arbitrary marked sets of one
orientation disseminate by positive Schatten H\"older through the two
transverse and two in-plane directions, with the original denominator at each
step. H\"older over six orientations and the six-witness, multiplicity-two
ledger give the ORIGINAL bound
\[
 \mathbb P(D\subset\mathcal B)\le p_B^{|D|},\quad
 p_B=6K_t^{1/12}e^{-tB(1-\cos\epsilon_0)/108}.                            \tag{1}
\]
It is unconditional, or averaged under the true initial-slice marginal, not
uniform in the fixed initial slice. This exact scope is used below.

\paragraph{1. Good-run spectral approximation in original coordinates.}
Take a consecutive good interval with $L+2$ transfer steps, $L\ge2$.
Transport it exactly to the small-bath frame. For its middle $L$ steps, replace
ONLY the two outermost half-onsite multipliers by the flat multiplier
$M_*(u)=e^{-2Bs(u)}$. Denote this auxiliary middle product by $A_L$.
This is an algebraic endpoint factorization: the original interval equals
$G=A_- A_L A_+$, with the removed positive endpoint ratios absorbed in $A_-$
and $A_+$. No integral or bath law has been replaced. The auxiliary product
closes with consistent flat endpoint data and satisfies the fixed-threshold
version of the path centered growing bath tubes calculation. Its exact Perron eigenvalue is $\lambda$, with positive
right/left vectors $v,l$, $\|v\|_2=1$, $l(v)=1$.

In path-centered coordinates the cone proof gives a complementary restriction
norm $\le\lambda r^L$, $r=2/5$, and a Perron projection norm $\le17/15$.
The full remainder includes $I-\Pi$, and is bounded by
$(32/15)\lambda r^L$. Returning to original coordinates multiplies this bound
by the norm of the exact exponential similarity. Its exponent is
$p\cdot x$, where $|p|\le3\sqrt B$ (four half-onsite gradients and one temporal
gradient, each unit angular derivative bounded by one) and
$|x|<\pi\sqrt B$. Thus its condition number is at most $e^{6\pi B}$:
\[
 \|A_L-\lambda v\otimes l\|\le(32/15)\lambda e^{6\pi B}r^L.             \tag{2}
\]
Haar translations are unitary and do not change this estimate.

\paragraph{2. Positivity upgrades(2) to a pointwise relative estimate.}
For any original chain product, its dependence on the first endpoint consists
of its first half-onsite multiplier (log oscillation at most $4B$) and one
temporal factor (at most $2B$). Its positive kernel therefore has endpoint
row ratio at most $e^{6B}$, independently of the number of factors. The same
holds at the last endpoint. The eigenfunction equations imply
\[
 \sup v/\inf v\le e^{6B},\quad \sup l/\inf l\le e^{6B},\quad
 \|v\|_1\ge e^{-6B},\quad \|l\|_1\ge e^{-6B}\|l\|_2\ge e^{-6B}.          \tag{3}
\]
All Haar measures are normalized. The endpoint ratios used in $A_\pm$ each
have log oscillation at most $8B$, so their relevant row/column oscillations
are at most $14B$. Hence
\[
 (A_-v)(x)\ge e^{-20B}\|A_-(x,\cdot)\|_2,
 \quad (lA_+)(y)\ge e^{-20B}\|A_+(\cdot,y)\|_2.                         \tag{4}
\]
Cauchy--Schwarz applied to(2), and then(4), yield
\[
 (1-\delta)\lambda a(x)b(y)\le G(x,y)
       \le(1+\delta)\lambda a(x)b(y),\quad
 \delta=3e^{64B}(2/5)^L,                                                \tag{5}
\]
where $a=A_-v$, $b=lA_+$ are positive. The constant $64$ safely exceeds
$6\pi+40$. Endpoint frame unitaries only change the positive functions and
preserve(5). This bound does not assert that $\lambda$ is the ground scalar
of the full period.

\paragraph{3. An arbitrary positive complement costs nothing further.}
Left or right multiplication by ANY positive integral kernels preserves the
relative rank-one inequalities(5). In particular the rest of a period may
contain arbitrarily many and arbitrarily bad original Wilson slices. Write
$W$ for the resulting whole-period operator and $t>0$ for the eigenvalue of
the resulting rank-one comparison kernel. Positivity gives
\[
 (1-\delta)t\le\rho(W)\le(1+\delta)t,
\quad
 (1-\delta)^q t^q\le\operatorname{Tr}W^q\le(1+\delta)^q t^q.             \tag{6}
\]
The first statement follows by testing the positive rank-one right function
in the eigenvalue comparison; the second follows by multiplying kernels
around the exact closed cycle. Thus, for $\delta\le1/2$,
\[
 e^{-4q\delta}\le\operatorname{Tr}W^q/\rho(W)^q\le e^{4q\delta}.          \tag{7}
\]
There is NO factor $e^{CB}$ per bad slice or per bad cluster. Only the single
good-run upgrade cost in(5) has been paid. The exact whole-period Perron scalar
is retained in(7); every raw scalar of the original complement is included.

\paragraph{4. A good run in each half-period controls the selected midpoint.}
For fixed bath, suppose each of the two length-$\ell$ arcs contains a good
run as above. Each full arc kernel then has relative rank-one error $\delta$,
regardless of its other slices. The conditional midpoint density with initial
selected link $x$ is proportional to the product of these two arc kernels.
For any two initial values $x,x'$, its projective density distance is at most
\[
 4\log\frac{1+\delta}{1-\delta}\le16\delta\quad(\delta\le1/2).
\]
Normalization constants cancel in projective distance. The elementary
Hellinger bound consequently gives
\[
 \operatorname{Aff}(p_x,p_{x'})^2\ge\operatorname{sech}^2(4\delta)
                  \ge1-16\delta^2.                                   \tag{8}
\]
This is uniform in initial values and retains actual arbitrary bath twists.
It controls selected conditional fibers, not the full bath affinity.

\paragraph{5. Price the run requirement with the proved original record.}
Partition an arc into $J=\lfloor\ell/(L+4)\rfloor$ disjoint blocks, each of
$L+4$ candidate good labels. The inner $L+2$ transfer steps then have all
required endpoint onsite and temporal labels. If none is entirely good, choose a bad label in
each block. The selected labels are distinct, so(1) and a union bound imply
\[
 \mathbb P\{\hbox{no good block in an arc}\}
       \le[(L+4)p_B]^J.                                                \tag{9}
\]
No independence is used. For $N_I$ plaquette-disjoint selected trajectories,
the probability that some trajectory lacks a good block in either arc is at
most $2N_I[(L+4)p_B]^J$. All such bounds remain averaged under the original initial-slice law.

For example choose
\[
 L=\max\{2,\lceil[64B+10\log(2+\ell)]/\log(5/2)\rceil\}.              \tag{10}
\]
Then $\delta\le3(2+\ell)^{-10}$. If $B\sim b\log\ell$ with ANY fixed
$b>0$, $K_t$ is uniformly bounded in its stated physical torus regime,
$L=O(B+\log\ell)$ and $(L+4)p_B\to0$. Equation(9) decays as $e^{-c\ell}$
for some fixed $c>0$ depending on $b,\epsilon_0,t$. This pays the deterministic
$e^{64B}$ cost by good-run length, rather than by a defective-slice exponential
moment. It permits many exceptional slices, whose total expected upper count
can diverge.

\paragraph{6. The remaining full-bath factor is explicit.}
Let $\eta_X$ be the ACTUAL bath law conditioned on the full initial slice,
including all selected-chain normalizers, and $p(X)$ its original marginal.
On the good-record event the selected conditional affinities multiply, giving
squared factor at least $a^2=(1-16\delta^2)^{N_I}$. Lift to the joint bath and
selected midpoint, then project. Writing
$\mathcal P_{\rm bath}=E_{X,Y\ {
m iid}\ p}\operatorname{Aff}(\eta_X,\eta_Y)^2$,
the original physical midpoint purity obeys the valid lower comparison
\[
 \mathcal P_{\rm full}\ge a^2\left[\mathcal P_{\rm bath}
                   -4N_I[(L+4)p_B]^J\right].                           \tag{11}
\]
Indeed the excluded overlap mass averaged in $X,Y$ is at most its unconditional
probability, and $(A-b)^2\ge A^2-2b$ for $0\le b\le A\le1$.
For polynomial $N_I$ in $1/a$, fixed physical arc length and
$B\sim b\log(1/a)$, both losses in(11) vanish. This extends conditional
microscopic decoupling beyond the earlier worst-bath $b<1/2$ restriction.
It does NOT bound $\mathcal P_{\rm bath}$ from below. Induced short-time spatial
interactions and the true remaining global modes remain in that law. No
spatial iteration, full S94 closure or mass gap is claimed.
\end{document}
